\subsection{Matrix Equivalence and Rank Normal Forms}
\begin{definition} \label{def:matrix-equivalence} \hlt{Matrix Equivalence}\\
For $A,B\in\operatorname{Mat}(m,n,\mathbb{F})$, write $A\sim_{\mathrm{eq}} B$ if $\exists$ $P\in\operatorname{Mat}(m,m,\mathbb{F})^{\times}$ and $Q\in\operatorname{Mat}(n,n,\mathbb{F})^{\times}$ s.t. $B=PAQ$.
\end{definition}

\begin{proposition} \label{prop:matrix-equivalence-basic-properties} \hlt{Matrix Equivalence and Elementary Operations}\\
The relation $\sim_{\mathrm{eq}}$ is an equivalence relation on $\operatorname{Mat}(m,n,\mathbb{F})$. Moreover, $A\sim_{\mathrm{eq}}B$ if and only if $B$ can be obtained from $A$ by finitely many elementary row and column operations.
\begin{proof}
Reflexivity follows from $A=I_mAI_n$. If $B=PAQ$ with $P,Q$ invertible, then $A=P^{-1}BQ^{-1}$, so symmetry holds. If $B=P_1AQ_1$ and $C=P_2BQ_2$, then $C=(P_2P_1)A(Q_1Q_2)$, and products of invertible matrices are invertible by Corollary \ref{cor:matrix-inverse-product-and-transpose}; hence transitivity holds.\\
Each elementary row operation is left multiplication by an invertible elementary matrix by Proposition \ref{prop:elementary-matrices-perform-row-operations}. Each elementary column operation is right multiplication by the transpose of an invertible elementary matrix, again by Proposition \ref{prop:elementary-matrices-perform-row-operations} and Corollary \ref{cor:matrix-inverse-product-and-transpose}. Thus any finite sequence of elementary row and column operations has the form $B=PAQ$ with $P,Q$ invertible, so $A\sim_{\mathrm{eq}}B$.\\
Conversely, suppose $A\sim_{\mathrm{eq}}B$, so $B=PAQ$ with $P,Q$ invertible. Lemma \ref{lem:invertible-matrices-row-reduce-to-identity} gives products of elementary matrices expressing $P$ and $Q$ as products of elementary factors. Left multiplication by the factors of $P$ is a sequence of elementary row operations, and right multiplication by the factors of $Q$ is a sequence of elementary column operations. Hence $B$ can be obtained from $A$ by finitely many elementary row and column operations.
\end{proof}
\end{proposition}

\begin{theorem} \label{thm:matrix-equivalence-same-linear-map} \hlt{Equivalent Matrices Represent the Same Linear Map}\\
Let $\dim\mathcal{V}=n$, $\dim\mathcal{W}=m$, and $A,B\in\operatorname{Mat}(m,n,\mathbb{F})$.\\
Then $A\sim_{\mathrm{eq}}B$ $\Longleftrightarrow$ $\exists\,T\in\LL(\mathcal{V},\mathcal{W})$ with $A=[T]_{\gamma\leftarrow\beta},\ B=[T]_{\gamma'\leftarrow\beta'}$ for some bases $\beta,\beta'$ of $\mathcal{V}$ and $\gamma,\gamma'$ of $\mathcal{W}$.
\begin{proof}
\begin{enumerate}[label=\roman*.]
\setlength{\itemsep}{0pt}
\item[$(\Rightarrow)$] Let $B=PAQ$ with $P,Q$ invertible. Choose bases $\beta,\gamma$ and define $T\in\LL(\mathcal{V},\mathcal{W})$ by $[T]_{\gamma\leftarrow\beta}=A$.\\
By Theorem \ref{thm:change-of-basis-matrices}, choose basis $\gamma'$ with $P=C_{\gamma'\leftarrow\gamma}$, and basis $\beta'$ with $Q^{-1}=C_{\beta'\leftarrow\beta}$. Then $Q=C_{\beta\leftarrow\beta'}$.\\
Then Theorem \ref{thm:two-sided-change-of-basis-for-linear-maps} gives $[T]_{\gamma'\leftarrow\beta'}=PAQ=B$.
\item[$(\Leftarrow)$] If $A=[T]_{\gamma\leftarrow\beta}$ and $B=[T]_{\gamma'\leftarrow\beta'}$, then $B=C_{\gamma'\leftarrow\gamma}AC_{\beta\leftarrow\beta'}$, so $A\sim_{\mathrm{eq}}B$.
\end{enumerate}
\hlt{Structural Commentary}: Equivalence is two-sided basis change.\\
\hlt{Key Insight}: Same map; different coordinates.
\end{proof}
\end{theorem}

\begin{theorem} \label{thm:rank-classifies-matrix-equivalence} \hlt{Rank Classifies Matrix Equivalence}\\
For $0\leq r\leq\min(m,n)$, set $J_r^{m,n}\defeq
\begin{bmatrix}
I_r & 0\\
0 & 0
\end{bmatrix}
\in\operatorname{Mat}(m,n,\mathbb{F})$.\\
For $A,B\in\operatorname{Mat}(m,n,\mathbb{F})$, $A\sim_{\mathrm{eq}}J_{\operatorname{rank}(A)}^{m,n}$ and $A\sim_{\mathrm{eq}}B\Longleftrightarrow \operatorname{rank}(A)=\operatorname{rank}(B)$.
\begin{proof}
\begin{enumerate}[label=\roman*.]
\setlength{\itemsep}{0pt}
\item[$(\Rightarrow)$] If $B=PAQ$ with $P,Q$ invertible, then $\operatorname{rank}(B)=\operatorname{rank}(A)$ by Theorem \ref{thm:full-rank-products-preserve-rank}.
\item[$(\Leftarrow)$] For any $X$ with $\operatorname{rank}(X)=r$, let $T_X:\mathbb{F}^n\to\mathbb{F}^m$ be $T_X(\mathbf{x})=X\mathbf{x}$.\\
Choose $\mathcal{M}_X$ with $\mathbb{F}^n=\ker(T_X)\oplus\mathcal{M}_X$; then $T_X|_{\mathcal{M}_X}:\mathcal{M}_X\to\operatorname{range}(T_X)$ is an isomorphism.\\
Choose a basis $(\mathbf{u}_1,\ldots,\mathbf{u}_r)$ of $\mathcal{M}_X$, extend it by a basis of $\ker(T_X)$ to $\beta_X$, and choose $\gamma_X$ beginning with $T_X\mathbf{u}_1,\ldots,T_X\mathbf{u}_r$. Then $[T_X]_{\gamma_X\leftarrow\beta_X}=J_r^{m,n}$, so $X\sim_{\mathrm{eq}}J_r^{m,n}$.\\
If $\operatorname{rank}(A)=\operatorname{rank}(B)=r$, then $A\sim_{\mathrm{eq}}J_r^{m,n}$ and $B\sim_{\mathrm{eq}}J_r^{m,n}$. By Proposition \ref{prop:matrix-equivalence-basic-properties}, symmetry gives $J_r^{m,n}\sim_{\mathrm{eq}}B$, and transitivity gives $A\sim_{\mathrm{eq}}B$.
\end{enumerate}
\hlt{Structural Commentary}: Rank is the invariant of two-sided equivalence.\\
\hlt{Key Insight}: Adapt bases to kernel and range.
\end{proof}
\end{theorem}

\begin{corollary} \label{cor:linear-map-rank-normal-form} \hlt{Rank Normal Form for Linear Maps}\\
Let $T\in\LL(\mathcal{V},\mathcal{W})$, $\dim\mathcal{V}=n$, $\dim\mathcal{W}=m$, and $\operatorname{rank}(T)=r$. Then $\exists$ bases $\beta,\gamma$ with $[T]_{\gamma\leftarrow\beta}=J_r^{m,n}$.
\begin{proof}
For initial bases $\beta_0,\gamma_0$, let $A=[T]_{\gamma_0\leftarrow\beta_0}$.\\
By Theorem \ref{thm:rank-classifies-matrix-equivalence}, $A\sim_{\mathrm{eq}}J_r^{m,n}$; Theorem \ref{thm:matrix-equivalence-same-linear-map} gives bases $\beta,\gamma$ with $[T]_{\gamma\leftarrow\beta}=J_r^{m,n}$.\\
\hlt{Structural Commentary}: Rank gives the two-basis normal form.\\
\hlt{Key Insight}: Rank normalizes the map.
\end{proof}
\end{corollary}

\begin{corollary} \label{cor:linear-map-equivalence-by-rank} \hlt{Linear Maps Classified by Rank}\\
For $S,T\in\LL(\mathcal{V},\mathcal{W})$, write $S\sim_{\mathrm{bc}}T$ if $\exists$ bases $\beta_S,\beta_T$ of $\mathcal{V}$ and $\gamma_S,\gamma_T$ of $\mathcal{W}$ such that $[S]_{\gamma_S\leftarrow\beta_S}=[T]_{\gamma_T\leftarrow\beta_T}$.\\
Then $S\sim_{\mathrm{bc}}T$ $\Longleftrightarrow$ $\operatorname{rank}(S)=\operatorname{rank}(T)$
\begin{proof}
\begin{enumerate}[label=\roman*.]
\setlength{\itemsep}{0pt}
\item[$(\Rightarrow)$] If $[S]_{\gamma_S\leftarrow\beta_S}=[T]_{\gamma_T\leftarrow\beta_T}$, then the two maps have the same matrix rank, hence the same rank.
\item[$(\Leftarrow)$] If $\operatorname{rank}(S)=\operatorname{rank}(T)=r$, Corollary \ref{cor:linear-map-rank-normal-form} gives bases with $[S]_{\gamma_S\leftarrow\beta_S}=J_r^{m,n}=[T]_{\gamma_T\leftarrow\beta_T}$.
\end{enumerate}
\hlt{Structural Commentary}: Rank is the invariant of basis-change equivalence.\\
\hlt{Key Insight}: Same rank means same normal form.
\end{proof}
\end{corollary}
